\PassOptionsToPackage{table,xcdraw}{xcolor}
\documentclass[11pt, a4paper]{article}

\usepackage[T1]{fontenc}
\usepackage[utf8]{inputenc}
\usepackage{tgpagella}
\usepackage[spanish, es-noshorthands]{babel}
\usepackage{amsmath, amssymb}
\usepackage[most]{tcolorbox}
\usepackage[margin=2.5cm]{geometry}
\usepackage{fancyhdr}

\pagestyle{fancy}
\fancyhf{}
\setlength{\headheight}{16pt}
\lhead{\textbf{UCB Sede Tarija} | Ing. Mecatrónica}
\rhead{IMT-221 | Soluciones S07-2}
\renewcommand{\headrulewidth}{0.4pt}

\begin{document}

\section*{Solucionario y Notas Docentes}

\subsection*{Diagnóstico Transversal (Esquema Colector Común)}
\begin{enumerate}
    \item \textbf{Identificación:} La señal entra por la \textbf{base}, sale por el \textbf{emisor}. El \textbf{colector} es el terminal común porque al apagar fuentes DC, $V_{CC}$ se vuelve Tierra AC[cite: 15].
    \item \textbf{Parámetros:} $g_m = I_{CQ}/V_T$ y $r_e \approx 1/g_m$ (o $r_\pi = \beta/g_m$).
    \item \textbf{Equivalente AC:} El alumno debe dibujar el modelo (Híbrido-$\pi$ o T), asegurando que el Colector aterrice en el símbolo de tierra y $R_E$ esté conectada desde emisor a tierra.
    \item \textbf{Leyes de Kirchhoff:} Si usa Modelo T: KVL en la malla de entrada $\implies v_i = i_e r_e + i_e R_E$.
    \item \textbf{Supuestos:} Banda media, operación en activa-directa, pequeña señal[cite: 15].
\end{enumerate}
\textbf{Error Común:} Memorizar que $A_v = -g_m R_C$ sin ver el circuito. \textbf{Sonda pedagógica:} "\textit{Mira por dónde sale la señal. ¿Hay un $R_C$ en el camino de salida en tu diagrama AC?}"[cite: 15].

\subsection*{Ejercicio Guiado: Par Diferencial DC}
\textbf{Paso 1: Estado Balanceado}
\begin{enumerate}
    \item KCL en nodo E: $I_{E1} + I_{E2} = I_{TAIL}$. Por simetría ($v_1 = v_2$), $I_{E1} = I_{E2}$. Sustituyendo: $2I_{E1} = 2\,\text{mA} \implies \mathbf{I_{E1} = 1\,\text{mA}}$ y $\mathbf{I_{E2} = 1\,\text{mA}}$[cite: 15].
    \item Asumiendo $\alpha \approx 1$, entonces $\mathbf{I_{C1} \approx 1\,\text{mA}}$ e $\mathbf{I_{C2} \approx 1\,\text{mA}}$[cite: 15].
    \item Ley de Ohm en colectores: $V_{RC1} = I_{C1}R_C = (1\,\text{mA})(4.7\,\text{k}\Omega) = \mathbf{4.7\,\text{V}}$. Igual para $V_{RC2} = 4.7\,\text{V}$[cite: 15].
    \item Potencial en colectores (KVL desde fuente): $V_{C1} = V_{CC} - V_{RC1} = 10\,\text{V} - 4.7\,\text{V} = \mathbf{5.3\,\text{V}}$. Igual para $V_{C2} = \mathbf{5.3\,\text{V}}$[cite: 15].
    \item \textbf{Verificación:} Si $V_B = 0\,\text{V}$ y la unión conduce, $V_E = V_B - V_{BE} = 0 - 0.7 = \mathbf{-0.7\,\text{V}}$. \\
    Voltaje Colector-Emisor: $V_{CE1} = V_{C1} - V_E = 5.3\,\text{V} - (-0.7\,\text{V}) = \mathbf{6.0\,\text{V}}$. \\
    Como $V_{CE} > 0.3\,\text{V}$ (aprox), el BJT no está saturado; opera correctamente en \textbf{Activa Directa}[cite: 15].
\end{enumerate}
\textbf{Error Común:} Decir que $V_C$ aumenta si la corriente aumenta. \textbf{Sonda pedagógica:} "\textit{Escribe la ecuación $V_C = V_{CC} - I_C R_C$. Si el término que restas se hace más grande, ¿qué le pasa a $V_C$?}"[cite: 15].

\textbf{Paso 2: Direccionamiento Cualitativo}
\begin{itemize}
    \item $\mathbf{v_1 > v_2}$: $I_{C1} \uparrow$, $I_{C2} \downarrow$. Consecuencia: $V_{C1} \downarrow$ (cae), $V_{C2} \uparrow$ (sube)[cite: 15].
    \item $\mathbf{v_1 < v_2}$: $I_{C1} \downarrow$, $I_{C2} \uparrow$. Consecuencia: $V_{C1} \uparrow$ (sube), $V_{C2} \downarrow$ (cae)[cite: 15].
\end{itemize}
\textit{Interpretación:} La naturaleza del Par Diferencial convierte un desbalance de voltaje en las bases en una redirección competitiva de la corriente de cola[cite: 15].

\end{document}
